The proposed framework rests on the premise that knowledge graph ontology is not a
neutral design choice but an active determinant of downstream predictive performance.
If this is true, then iteratively adapting the ontology in response to task feedback
should yield measurable improvements --- both in accuracy and in the structural
properties of the resulting graph. We formalise this premise through three testable
hypotheses, each targeting a distinct dimension of the expected benefit: predictive
accuracy, graph efficiency, and structural informativeness.

\begin{hypothesis}\label{hyp:accuracy}
Task-aware knowledge graph construction improves prediction accuracy in node-level
price movement classification compared to static graph construction approaches.
\end{hypothesis}

Ontology configurations that are better aligned with the prediction objective should
produce graph-derived features with higher discriminative power. H\ref{hyp:accuracy}
is evaluated by comparing AUC and F1-score of the downstream classifier trained on
features extracted from task-aware KG variants against a static baseline in which the
ontology is fixed throughout (Section~\ref{sec:results}).

\begin{hypothesis}\label{hyp:efficiency}
Task-aware optimization enables smaller knowledge graphs while maintaining or
improving downstream task performance.
\end{hypothesis}

By discarding ontology elements that do not contribute to predictive performance, the
feedback loop should converge on more compact representations. H\ref{hyp:efficiency}
is evaluated by tracking the number of nodes and edges per iteration alongside
classifier AUC, testing whether the task-aware variant achieves a favourable
size--performance tradeoff relative to the unconstrained baseline
(Section~\ref{sec:results}).

\begin{hypothesis}\label{hyp:structure}
Iterative, task-driven knowledge graph construction produces more informative graph
structures and facilitates the identification of effective ontology configurations.
\end{hypothesis}

Successive ontology revisions guided by performance feedback should increase the
diversity and predictive relevance of extracted relation types. H\ref{hyp:structure}
is evaluated by examining the ontology evolution trace --- which relation types are
added, retained, or removed across iterations --- and by measuring feature entropy
and metapath diversity over successive construction steps
(Section~\ref{sec:results}).
